% !TEX root = standalone.tex
% Research gap analysis (issue #52).
% Cross-references into the two foundations use \xsec{<standalone number>}{<label>}:
% the standalone build prints the number, thesis-draft.tex prints a live \ref.
\providecommand{\unitname}{chapter}
\providecommand{\xsec}[2]{#1}
\providecommand{\aimref}{the following chapter}

\chapter{Research Gap}
\label{chap:gap}

The two foundations answer different questions. The biological foundation
describes the circuit the model follows: a laminar column whose inhibition is
divided among interneuron classes that act on different compartments of the
pyramidal cell. The computational foundation describes what existing models do
with sparse coding, interneuron classes and associative memory, and places them
on three axes: substrate, inhibition model, and the arbitration between
encoding and retrieval. This \unitname{} draws the two together. It first sets
out what is already established and is therefore not claimed here
(Section~\ref{sec:gap-established}), then shows where these strands fall
short and states the gap (Section~\ref{sec:gap-open}), and finally answers the objections that a reader
familiar with the field is most likely to raise
(Section~\ref{sec:gap-objections}).

\section{What Is Established}
\label{sec:gap-established}

Most of the components of this thesis are not new, and the claim made here
depends on keeping them apart from what is.

The compartmental division of inhibition is established biology. PV basket and
chandelier cells inhibit the perisomatic region and the axon initial segment,
SST Martinotti cells inhibit the apical tuft, and VIP cells disinhibit the tuft
by inhibiting SST cells~\cite{tremblay2016}
(Section~\xsec{1.3}{sec:bio-inhibitory}). That PV and SST cells play different
roles in a circuit is also established in models. Rate models with
multi-compartment pyramidal cells have shown how the VIP, SST and PV motif
gates pathways and redistributes inhibition between soma and
dendrite~\cite{yang2016gating,hertag2019amplifying}, and laminar spiking models
now place PV, SST and VIP populations in every layer of a
column~\cite{billeh2020systematic,moreni2025celltype}.

The two operators that the thesis assigns to these classes are established as
well. Divisive normalisation is a canonical neural
computation~\cite{carandini2012normalization}, and its link to PV cells has
direct, if contested, experimental support~\cite{wilson2012division}.
Dendritic gating of top-down input at the apical tuft is established at the
level of single cells and circuits~\cite{larkum1999new,larkum2013,murayama2009dendritic},
and has been turned into a learning principle in which SST-like interneurons
act on the apical
compartment~\cite{urbanczik2014learning,sacramento2018dendritic}.

The use of interneuron classes inside a memory network is not new either. In
the attractor and BCPNN tradition, basket cells implement competition between
minicolumns, Martinotti cells set how long an attractor lasts, and double
bouquet cells make learned inhibition compatible with Dale's
principle~\cite{krishnamurthy2012cortical,chrysanthidis2019introducing}.
Finally, the substrate is taken over unchanged from the HTM tradition: sparse
distributed representations over minicolumns with dendritic
segments~\cite{ahmad2016sdr,hawkins2016neurons,cui2017sp}, and the separation of
encoding and retrieval as a functional requirement of memory has a clear
precedent in the hippocampal model of Hasselmo et al.~\cite{hasselmo2002theta}.
Cutsuridis et al.~\cite{cutsuridis2010encoding} have shown, in a model of
hippocampal CA1, how PV-type and SST-type interneurons can divide the work of
the two modes between them.

\section{Where the Strands Fall Short}
\label{sec:gap-open}

Each of these strands, then, is well developed on its own terms. What they
have not done is come together. Table~\ref{tab:gap-matrix} shows this by
reducing Table~\xsec{2.1}{tab:comp-prior-work} to a single row per tradition,
set against the three parts of the gap. No tradition covers more than two of
them. The one model that does test how interneuron classes share encoding and
retrieval is a biophysical model of hippocampal CA1, not a neocortical SDR
column.

\begin{table}[htbp]
\centering
\small
\setlength{\tabcolsep}{3pt}
\renewcommand{\arraystretch}{1.15}
\begin{tabular}{>{\raggedright\arraybackslash}p{2.9cm}>{\raggedright\arraybackslash}p{2.8cm}>{\raggedright\arraybackslash}p{3.0cm}>{\raggedright\arraybackslash}p{3.0cm}}
\hline
\textbf{Tradition} & \textbf{(i) Discrete SDR substrate} & \textbf{(ii) Operator-distinct PV and SST} & \textbf{(iii) Differential PV/SST roles in encoding and retrieval tested} \\
\hline
Hippocampal CA1 model & No: biophysical spiking cells & Partly: PV-type and SST-type cells, not formal operators & Yes, with the mode set by septal theta \\
HTM, CLA and SDR models & Yes & No: homogeneous kWTA & No: single inhibition rule \\
SDM and NAA & Yes (NAA); abstract binary (SDM) & No & No: single inhibition rule \\
Interneuron-resolved circuit models & No: rate or spiking populations & Yes & No: no memory task \\
Associative-memory columns & Partly: minicolumns, no learned SDRs & Partly: several classes, roles in retrieval dynamics & No: storage and recall roles not compared \\
Operator mechanisms & No & Partly: operators defined, cell attribution contested & No: no memory task \\
\hline
\textbf{This thesis} & \textbf{Yes} & \textbf{Yes} & \textbf{Tested, with the mode set externally} \\
\hline
\end{tabular}
\caption{The research gap at a glance. Each row condenses one tradition from
Table~\xsec{2.1}{tab:comp-prior-work}; the columns are the three parts of the
gap stated in the text.}
\label{tab:gap-matrix}
\end{table}

Put together, the gap has three parts: (i) a discrete SDR substrate in the HTM
tradition; (ii) inhibition split into two interneuron types with distinct
operations, a PV-like divisive normalisation on the output compartment and an
SST-like dendritic gate on the apical input compartment; and (iii) an explicit
test of whether these two operations contribute differently to encoding and to
retrieval. This thesis sets out to hold all three at once. In practice, that
means taking a biological distinction that is already well established,
carrying it into a substrate that has so far reduced inhibition to a single
sorting rule, and asking whether the move changes anything that homogeneous
kWTA could not. As in the hippocampal models, which mode is active is set from
outside the column, by the encoding or retrieval condition presented to it.
What the thesis tests is how the two interneuron types share the work within
each mode, and, where feedforward and recalled signals arrive together, how the
column integrates them.

The contribution is therefore a modest one. The thesis does not offer a new
account of cortical inhibition or a new memory model, but a controlled
substitution inside an existing one, and the outcome is useful either way. If
the two operations do divide encoding and retrieval between them, the SDR
substrate gains a cell-type-specific division of labour between encoding and
retrieval, of the kind Cutsuridis et al. describe for CA1. If they do not, that is evidence that, at least in this setting, the
homogeneous rule loses less than the biology would suggest.

\section{Anticipated Objections}
\label{sec:gap-objections}

\paragraph{This has already been done in the hippocampus.}
Cutsuridis et al.~\cite{cutsuridis2010encoding} assign storage and recall to
different interneuron classes in CA1 and test those roles by removing the
classes one at a time. This is the closest precedent for the thesis, and the
thesis follows it in letting the mode be set from outside the circuit. The
model, however, is a compartmental spiking model of a hippocampal region whose
inputs arrive in separate dendritic layers from entorhinal cortex and CA3, and
its interneurons are detailed cell models rather than operators that could
replace a selection rule. The question of what happens when the single
inhibition rule of an SDR column is replaced by two such operators, in a
neocortical laminar setting, does not arise there.

\paragraph{PV and SST have already been dissociated.}
Hertäg and Sprekeler~\cite{hertag2019amplifying} model PV cells on the soma and
SST cells on the apical dendrite, and show that the circuit can switch
persistently between integrating and cancelling top-down input. This is the
closest precedent for a circuit-internal switch and the thesis treats it as
such. Its substrate, however, is a rate model with no sparse distributed
representations and no learned associations, and the switch concerns how
feedback is processed, not whether a pattern is stored or recalled.

\paragraph{Interneuron classes have already been studied in memory networks.}
Krishnamurthy et al.~\cite{krishnamurthy2012cortical,krishnamurthy2015longrange}
show that basket and Martinotti cells shape attractor memory differently, with
Martinotti cells setting dwell time and producing surround suppression. Here
the memory setting is real, but the roles concern the dynamics of a pattern
that is already being recalled, in a spiking attractor network whose patterns
were not learned as SDRs, and the change of mode in this tradition comes from
neuromodulation~\cite{fransen1998horizontal}.

\paragraph{The mechanisms are borrowed rather than discovered.}
This is correct and intended. The thesis does not propose a new form of
inhibition. Its claim concerns the consequence of placing two known operators
in a substrate that lacks them, which is a question the source models cannot
answer because they do not use that substrate. A negative result, in which the
two operators behave no differently from kWTA with respect to encoding and
retrieval, would be as informative as a positive one.

\paragraph{PV is not purely divisive and SST is not purely subtractive.}
The arithmetic of each class depends on stimulus and network
state~\cite{elboustani2014response,seybold2015}, and the attribution
of division to PV cells is contested~\cite{lee2012activation,atallah2012parvalbumin}.
The thesis therefore makes its claim about operators, not about cell classes.
The operators are called PV-like and SST-like because their targets follow the
compartmental targeting described in Section~\xsec{1.3}{sec:bio-inhibitory},
which is robust, while the mapping of each operator to a single class is a
modelling assumption on which the results about the operators do not depend.

\paragraph{The HTM tradition already names its inhibitory cells.}
Hawkins and Ahmad~\cite{hawkins2016neurons} identify fast-spiking basket cells as
the likely source of the inhibition their model needs, and the author of NAA
leaves the inhibition model explicitly open~\cite{dobric2024naa}. Naming a cell
type is not implementing it: in both cases the computation remains a
homogeneous k-winners-take-all rule
(Section~\xsec{2.2.5}{subsec:comp-htm-inhibition}).

The aim of the thesis follows directly from this gap and is stated in
\aimref{}.
